% Generated draft: requires mathematical and source review.
\documentclass[11pt]{article}
\usepackage[margin=25mm]{geometry}
\usepackage{fontspec}
\setmainfont{Latin Modern Roman}
\usepackage{amsmath,amssymb,mathtools,unicode-math}
\setmathfont{Latin Modern Math}
\usepackage{enumitem,needspace,xurl,hyperref}
\setlength{\parindent}{0pt}
\setlength{\parskip}{5pt}
\setlength{\emergencystretch}{3em}
\newcommand{\N}{\mathbb{N}}
\newcommand{\Z}{\mathbb{Z}}
\newcommand{\Q}{\mathbb{Q}}
\newcommand{\R}{\mathbb{R}}
\newcommand{\C}{\mathbb{C}}
\newcommand{\F}{\mathbb{F}}
\newcommand{\A}{\mathbb{A}}
\newcommand{\PP}{\mathbb{P}}
\newcommand{\E}{\mathbb{E}}
\newcommand{\eps}{\varepsilon}
\newcommand{\dd}{\,\mathrm{d}}
\begin{document}
\section*{UnsolvedMath ID: 11000251}
\textbf{Source problem code:} catalogue entry
\subsection{Definitions and mathematical statement}
Let \(\Sigma_2\) be a closed connected orientable surface of genus two. The question asks whether there exists a closed connected four-manifold \(M\), carrying a complete real hyperbolic metric of constant sectional curvature \(-1\), together with a locally trivial surface bundle
\[\Sigma_2\longrightarrow M\xrightarrow{p}\Sigma_2.\]
Thus both base and fiber have genus exactly two. In the usual smooth interpretation, the bundle has local trivializations \(p^{-1}(U)\cong U\times\Sigma_2\). The hyperbolic metric need not make the bundle projection a Riemannian submersion or the fibers totally geodesic; neither requirement is in the source. ``Hyperbolic four-manifold'' here has its ordinary real-hyperbolic meaning, not complex hyperbolicity or just word hyperbolicity of a group.

The source does not explicitly repeat smoothness/orientation conventions or require that monodromy preserve a chosen fiber orientation. If orientation-preserving structure group is intended, that must be retained from context rather than inferred only from the abstract orientability of the fiber.

The following paragraph begins Section 4.2 and reports evidence in [21] for a negative answer to \emph{Question 4.1}, then announces a conjecture that is a special case of a more general Seiberg--Witten-vanishing conjecture, referring to [37,14] for definitions. Question 4.1 and that conjecture are absent. Accordingly the reported negative evidence is not silently reassigned to a proved obstruction for this exact genus-two bundle question.
\subsection{Short English statement}
Does a closed real-hyperbolic four-manifold admit a surface-bundle structure with both fiber and base of genus two? The adjacent Seiberg--Witten paragraph refers to another question and does not supply an obstruction theorem here.
\subsection{Sources}
\begin{itemize}
\item Local UnsolvedMath registry, permanent ID \(11000251\), exact source code \texttt{catalogue entry}.
\item Farb, \emph{Problems on Mapping Class Groups and Related Topics}, Question 4.3, PDF page 288, followed by Section 4.2: \url{https://www.math.uchicago.edu/~farb/papers/mcgbook.pdf}. Registry access 2026-07-29; pdf-font-structure extraction.
\item Internal references [21] (evidence/conjecture) and [37,14] (definitions) lack full entries, and Question 4.1/the following conjecture are not supplied.
\end{itemize}
\end{document}
